% answers_short.tex —— 2004 年·春季上海卷 · 简短答案速查
% 由 tools/make_short_answers.py 自 answers.tex 的 \gksolhead/\gkanswer 提取，
% 并与 2004/2004_数学真题_春季上海_答案要点.md 逐题比对后生成，勿手改（改源 MD 后重跑）。

\begin{gknotice}
\textbf{使用说明}：本节只给最终答案便于快速核对，完整推导见解析版。答案为编写者独立推导，\textbf{非官方标准答案}。
\end{gknotice}

\begin{gknotice}
选择题答案中的 U 表示待核，未给出猜测答案，详情见解析版。
\end{gknotice}

\gkpart{填空题}
\begin{itemize}[leftmargin=2.6em]
  \item[\textbf{1.}] $1$
  \item[\textbf{2.}] $2$
  \item[\textbf{3.}] $2$
  \item[\textbf{4.}] $(x-1)^2+y^2=4$
  \item[\textbf{5.}] $1$
  \item[\textbf{6.}] $\arctan\dfrac14$
  \item[\textbf{7.}] $3$
  \item[\textbf{8.}] $n^2-n+1$
  \item[\textbf{9.}] $\dfrac5{14}$
  \item[\textbf{10.}] $4(x-3)^2+9(y-2)^2=36$
  \item[\textbf{11.}] 第 $34$ 行（以第 0 行为首行）
  \item[\textbf{12.}] 例如 $a_n=(-1)^{n-1}$
\end{itemize}

\gkpart{选择题}
\begin{center}\begin{tabular}{ccccc}
\toprule
题号 & 13 & 14 & 15 & 16 \\
\midrule
答案 & D & U & C & B \\
\bottomrule
\end{tabular}\end{center}

\gkpart{解答题}
\begin{itemize}[leftmargin=2.6em]
  \item[\textbf{17.}] $x=\dfrac\pi3$ 或 $x=\dfrac\pi2$
  \item[\textbf{18.}] 原文中约束不等式含 $x$ 而非 $p$，因此不能确定 $p$。按字面，方程有实根当且仅当 $|p|\ge2$；若原意为 $\dfrac{2p+1}{p+2}<0$，则 $-2<p< -\dfrac12$，方程无实根。此结论待原题变量疑点裁定。
  \item[\textbf{19.}] (1) $1458$ 辆；(2) 2011 年底
  \item[\textbf{20.}] (1) $CC_1\perp MN$；(2) $S_{AB}^2=S_{BC}^2+S_{AC}^2-2S_{BC}S_{AC}\cos\theta$
  \item[\textbf{21.}] (1) $a=1$；(2) $[-\dfrac12,+\infty)$；(3) 命题成立
  \item[\textbf{22.}] (1) $B=(4,1)$；(2) $a=2$；(3) $h(t)=\begin{cases}\sqrt{(t-1)^2+4},&t<-1,\\\dfrac{|t-3|}{\sqrt2},&-1\le t\le5,\\\sqrt{(t-4)^2+1},&t>5.\end{cases}$
\end{itemize}
